\section{Symbolregister}
\label{sec:symbolregister}

\small
\begin{longtable}{@{}p{.13\textwidth}p{.37\textwidth}p{.14\textwidth}p{.28\textwidth}@{}}
\toprule
Symbol & Betydning & Enhed/type & Dukker op i \\
\midrule
\endfirsthead
\toprule
Symbol & Betydning & Enhed/type & Dukker op i \\
\midrule
\endhead
\midrule
\multicolumn{4}{r}{Fortsættes på næste side}\\
\endfoot
\bottomrule
\endlastfoot
\multicolumn{4}{@{}l}{\textbf{Signaler og signalveje}}\\
\(r\), \(R(s)\), \(R\) & Reference i tid/Laplace-domæne. \(R(s)=1/s\) for unit step og \(R(s)=h_0/s\) for step med størrelse \(h_0\). & Systemspecifik & \S\ref{sec:start}, \S\ref{sec:feedback}, \S\ref{sec:disturbance}, \S\ref{sec:metodefinder}\\
\(e\), \(E(s)\), \(E\) & Fejl mellem reference og målt output; fx \(e=r-Hy\). & Systemspecifik & \S\ref{sec:start}, \S\ref{sec:feedback}, \S\ref{sec:disturbance}\\
\(u\), \(U(s)\), \(U\) & Kontrolsignal/plantinput. I modelafsnit er \(U(s)\) Laplace-transformeret input. & Systemspecifik & \S\ref{sec:model}, \S\ref{sec:feedback}, \S\ref{sec:limits}, \S\ref{sec:metodefinder}\\
\(u_{\mathrm{lin}}(t)\) & Ubegrænset lineær controllerkommando før saturation/rate limit. & Som \(u\) & \S\ref{sec:limits}, \S\ref{sec:python}\\
\(u_{\mathrm{sat}}(t)\) & Satureret aktuatorsignal efter amplitudebegrænsning. & Som \(u\) & \S\ref{sec:limits}\\
\(u_{\min},u_{\max}\) & Nedre og øvre amplitudelimit for aktuator. & Som \(u\) & \S\ref{sec:limits}\\
\(\dot u(t)\) & Tidsafledt aktuatorsignal; bruges til rate limit. & \(u\)/s & \S\ref{sec:limits}\\
\(y\), \(Y(s)\), \(Y\) & Output i tid/Laplace-domæne. I modelafsnit er \(Y(s)\) Laplace-transformeret output. & Systemspecifik & \S\ref{sec:start}, \S\ref{sec:model}, \S\ref{sec:feedback}, \S\ref{sec:disturbance}\\
\(d\), \(D(s)\), \(D\) & Disturbance og disturbance path. \(D(s)\) kan også være transferfunktionen fra disturbance til outputbidrag. & Varierer & \S\ref{sec:disturbance}, \S\ref{sec:feedforward}\\
\(d_0\) & Konstant disturbance-størrelse i \(D(s)=d_0/s\). & Som disturbance & \S\ref{sec:disturbance}\\
\(n\), \(N\) & Measurement noise i feedbackmålingen; \(N=0\) betyder at noise sættes til nul i superposition. & Varierer & \S\ref{sec:disturbance}\\
\(x\), \(x_0\), \(\tilde x\) & Tilstands-/positionsvariabel, arbejdspunkt og deviation variable ved linearisering. & Varierer & \S\ref{sec:model}\\
\(u_0\), \(\tilde u\) & Input-arbejdspunkt og input-deviation ved linearisering. & Som \(u\) & \S\ref{sec:model}\\
\(h_0\) & Stepstørrelse/amplitude. & Som reference & \S\ref{sec:time}, \S\ref{sec:feedback}\\
\multicolumn{4}{@{}l}{\textbf{Transferfunktioner og blokke}}\\
\(G(s)\), \(G\) & Plant eller generel transferfunktion. Bruges også i \(G=Y/U\). & Ratio & \S\ref{sec:start}, \S\ref{sec:metodefinder}, \S\ref{sec:model}, \S\ref{sec:bode}, \S\ref{sec:feedback}, \S\ref{sec:pilead}, \S\ref{sec:disturbance}, \S\ref{sec:feedforward}, \S\ref{sec:python}\\
\(G_1(s)\) & Førsteordens standardmodel. & Ratio & \S\ref{sec:time}\\
\(G_2(s)\) & Andenordens standardmodel. & Ratio & \S\ref{sec:time}\\
\(G_{er}(s)\) & Error transferfunktion \(E(s)/R(s)\). & Ratio & \S\ref{sec:feedback}\\
\(G_{yd}(s)\) & Transferfunktion fra disturbance til output. & Ratio & \S\ref{sec:feedforward}\\
\(F_{\mathrm{track}}(s)\) & Reference feed-forward, typisk plantinvers plus lavpasfilter. & Ratio & \S\ref{sec:tracking-feedforward}\\
\(C(s)\), \(C\) & Controller. & Ratio & \S\ref{sec:metodefinder}, \S\ref{sec:feedback}, \S\ref{sec:pilead}, \S\ref{sec:limits}, \S\ref{sec:feedforward}\\
\(C_0(s)\), \(C_0\) & Controllerdelen uden \(K_P\), brugt når \(K_P\) isoleres fra magnitudebalancen. & Ratio & \S\ref{sec:pilead}\\
\(C_{PI}\) & PI-controllerform/genkendelsessignal. & Ratio & \S\ref{sec:metodefinder}\\
\(C_{\mathrm{lead}}\) & Lead-controllerform/genkendelsessignal. & Ratio & \S\ref{sec:metodefinder}\\
\(C_{\mathrm{lag}}(s)\) & Lag-led i P-Lead-Lag-design. & Ratio & \S\ref{sec:limits}\\
\(H(s)\), \(H\) & Sensor-/måleled i feedbackgrenen. & Ratio & \S\ref{sec:feedback}, \S\ref{sec:disturbance}\\
\(F(s)\), \(F\) & Referenceprefilter. & Ratio & \S\ref{sec:disturbance}\\
\(F_d(s)\), \(F_d\) & Disturbance feed-forward-blok. & Ratio & \S\ref{sec:metodefinder}, \S\ref{sec:disturbance}, \S\ref{sec:feedforward}, \S\ref{sec:python}\\
\(L(s)\), \(L\) & Loop transfer function, typisk \(CG\) eller \(CGH\) afhængigt af sensor. & Ratio & \S\ref{sec:bode}, \S\ref{sec:feedback}, \S\ref{sec:disturbance}, \S\ref{sec:python}\\
\(S(s)\), \(S\) & Sensitivity \(1/(1+L)\). & Ratio & \S\ref{sec:metodefinder}, \S\ref{sec:feedback}, \S\ref{sec:disturbance}\\
\(T(s)\), \(T\) & Complementary sensitivity \(L/(1+L)\). & Ratio & \S\ref{sec:metodefinder}, \S\ref{sec:disturbance}\\
\multicolumn{4}{@{}l}{\textbf{Laplace, polynomier og stabilitet}}\\
\(s\) & Laplacevariabel. Bruges også som polynomiumsvariabel i scripts. & \(\mathrm{s}^{-1}\) & \S\ref{sec:model}, \S\ref{sec:feedback}, \S\ref{sec:pilead}, \S\ref{sec:limits}, \S\ref{sec:feedforward}, \S\ref{sec:python}\\
\(j\), \(\mathrm{j}\) & Imaginær enhed i frekvensrespons \(s=j\omega\). & -- & \S\ref{sec:model}, \S\ref{sec:bode}, \S\ref{sec:pilead}\\
\(\sigma\) & Realdel af \(s=\sigma+j\omega\). & \(\mathrm{s}^{-1}\) & \S\ref{sec:symbolregister}\\
\(\omega\) & Vinkelfrekvens. & rad/s & \S\ref{sec:start}, \S\ref{sec:model}, \S\ref{sec:bode}, \S\ref{sec:pilead}, \S\ref{sec:python}\\
\(p\) & Pol eller karakteristisk polynomium afhængigt af kontekst; \(p(s,K)\) er karakteristisk polynomium med gain. & \(\mathrm{s}^{-1}\) eller polynomium & \S\ref{sec:model}, \S\ref{sec:python}\\
\(z\) & Zero/nulpunkt. & \(\mathrm{s}^{-1}\) & \S\ref{sec:model}\\
\(a_n,\ldots,a_0\) & Nævnerkoefficienter i transferfunktion. \(a_1,a_0\) bruges også ved andenordensmatch. & Varierer & \S\ref{sec:model}, \S\ref{sec:time}\\
\(b_m,\ldots,b_0\) & Tællerkoefficienter i transferfunktion. & Varierer & \S\ref{sec:model}\\
\(m,n\) & Polynomiumsgrader eller afledningsorden. \(n\) kan også betyde noise; læs konteksten. & Heltal & \S\ref{sec:model}, \S\ref{sec:disturbance}\\
\(K\) & Generelt realt gain eller parameter i karakteristisk polynomium. & Varierer & \S\ref{sec:model}, \S\ref{sec:python}\\
\(K_{\mathrm{ss}}\) & DC/static gain \(\lim_{s\to0}G(s)\). & output/input & \S\ref{sec:model}, \S\ref{sec:time}\\
\(\LHP\) & Left half-plane; poler med negativ realdel. & Område & \S\ref{sec:start}, \S\ref{sec:model}, \S\ref{sec:feedback}, \S\ref{sec:feedforward}, \S\ref{sec:python}\\
\(\RHP\) & Right half-plane; poler med positiv realdel. & Område & \S\ref{sec:model}, \S\ref{sec:bode}, \S\ref{sec:nyquist}, \S\ref{sec:python}\\
\(P\) & Antal open-loop RHP-poler i Nyquistkriteriet. & Antal & \S\ref{sec:metodefinder}, \S\ref{sec:nyquist}\\
\(N\) & Netto clockwise encirclements i Nyquistkriteriet. Kan også betyde noise; læs konteksten. & Antal eller signal & \S\ref{sec:nyquist}, \S\ref{sec:disturbance}\\
\(Z\) & Antal closed-loop RHP-poler i Nyquistkriteriet. & Antal & \S\ref{sec:nyquist}\\
\multicolumn{4}{@{}l}{\textbf{Tidsrespons}}\\
\(\tau\) & Førsteordens tidskonstant. & s & \S\ref{sec:time}\\
\(\omega_b\) & Break frequency for førsteordensled. & rad/s & \S\ref{sec:time}\\
\(\omega_n\) & Naturlig frekvens for andenordensstandardform. & rad/s & \S\ref{sec:time}, \S\ref{sec:python}\\
\(\zeta\) & Damping ratio. & -- & \S\ref{sec:time}, \S\ref{sec:python}\\
\(M_p\) & Overshoot som ratio; procent fås ved at gange med 100. & -- eller \% & \S\ref{sec:time}, \S\ref{sec:python}\\
\(t_r\) & Rise time. & s & \S\ref{sec:metodefinder}, \S\ref{sec:time}\\
\(t_s\) & Settling time, her typisk 2\%-approksimation. & s & \S\ref{sec:metodefinder}, \S\ref{sec:time}\\
\multicolumn{4}{@{}l}{\textbf{Bode, Nyquist og marginer}}\\
\(M\) & Lineær magnitude, fx \(M=|G(j\omega_c)|\). & Ratio & \S\ref{sec:bode}, \S\ref{sec:pilead}\\
\(M_{\mathrm{dB}}\) & Magnitude i decibel. & dB & \S\ref{sec:bode}\\
\(\omega_c\) & Gain crossover frequency, hvor loopmagnitude er 1. & rad/s & \S\ref{sec:metodefinder}, \S\ref{sec:bode}, \S\ref{sec:pilead}, \S\ref{sec:limits}, \S\ref{sec:python}\\
\(\omega_\pi\) & Phase crossover frequency, hvor fasen er \(-180^\circ\). & rad/s & \S\ref{sec:bode}\\
\(\gamma_M\) & Phase margin. & grader & \S\ref{sec:metodefinder}, \S\ref{sec:bode}, \S\ref{sec:pilead}, \S\ref{sec:limits}, \S\ref{sec:python}\\
\(K_M\) & Gain margin som lineær faktor. & faktor & \S\ref{sec:bode}\\
\(\phi_c\) & Fase ved crossoverpunkt på enhedscirklen. & grader & \S\ref{sec:bode}\\
\(x,y\) & Real- og imaginærdel af Nyquistpunkt; også generelle tilstands-/outputsymboler i andre kontekster. & Ratio & \S\ref{sec:model}, \S\ref{sec:bode}\\
\(a\) & Positiv realakseskæringsstørrelse når Nyquistkurven skærer i \(-a\). & Ratio & \S\ref{sec:nyquist}\\
\multicolumn{4}{@{}l}{\textbf{Controllerdesign}}\\
\(K_P\) & Proportionalgain. & -- eller ratio & \S\ref{sec:start}, \S\ref{sec:metodefinder}, \S\ref{sec:bode}, \S\ref{sec:feedback}, \S\ref{sec:pilead}, \S\ref{sec:python}\\
\(K_u\) & Ultimativ/kritisk proportionalgain ved sustained oscillation. & -- eller ratio & \S\ref{sec:ziegler}\\
\(P_u\) & Ultimativ oscillationsperiode. & s & \S\ref{sec:ziegler}\\
\(\omega_u\) & Ultimativ oscillationsfrekvens. & rad/s & \S\ref{sec:ziegler}\\
\(\tau_i\) & PI-/Lag-tid; \(N_i=\omega_c\tau_i\). & s & \S\ref{sec:pilead}, \S\ref{sec:limits}, \S\ref{sec:python}\\
\(T_i\) & Integral time i ideal PID/Ziegler--Nichols notation. & s & \S\ref{sec:ziegler}\\
\(T_d\) & Derivative time i ideal PID/Ziegler--Nichols notation. & s & \S\ref{sec:ziegler}\\
\(\tau_d\) & Lead-tid. & s & \S\ref{sec:pilead}, \S\ref{sec:python}\\
\(\alpha\) & Lead-forhold; skal opfylde \(0<\alpha<1\). & -- & \S\ref{sec:start}, \S\ref{sec:metodefinder}, \S\ref{sec:pilead}, \S\ref{sec:limits}, \S\ref{sec:python}\\
\(N_i\) & Dimensionsløs PI-zero-placering \(N_i=\omega_c\tau_i\). & -- & \S\ref{sec:metodefinder}, \S\ref{sec:pilead}, \S\ref{sec:limits}, \S\ref{sec:python}\\
\(\phi_{PI}\) & PI-leddets fasebidrag ved crossover. & grader & \S\ref{sec:pilead}\\
\(\phi_G\) & Plantens fase ved crossover. & grader & \S\ref{sec:pilead}\\
\(\phi_{\mathrm{lead}}\) & Krævet leadfasebidrag. & grader & \S\ref{sec:pilead}\\
\(\phi_{\max}\) & Maksimalt leadfaseboost. & grader & \S\ref{sec:pilead}\\
\(\beta\) & Lag-forhold; skal opfylde \(\beta>1\). & -- & \S\ref{sec:start}, \S\ref{sec:metodefinder}, \S\ref{sec:limits}, \S\ref{sec:python}\\
\(\phi_{\mathrm{lag}}\) & Lag-leddets fasebidrag; normalt negativt. & grader & \S\ref{sec:limits}\\
\(\phi_{\mathrm{uden\ lag}}\) & Samlet fase før Lag-leddet tilføjes. & grader & \S\ref{sec:limits}\\
\multicolumn{4}{@{}l}{\textbf{Limits, feed-forward og øvrige}}\\
\(R_{\max}\) & Maksimal aktuatorrate. & \(u\)/s & \S\ref{sec:limits}\\
\(\tau_f\) & Lavpasfilterets tidskonstant i en realiserbar feed-forward invers. & s & \S\ref{sec:tracking-feedforward}\\
\(\sigma_D\) & Fortegn for disturbancebidrag ved summering; \(\sigma_D\in\{-1,+1\}\). & -- & \S\ref{sec:feedforward}, \S\ref{sec:python}\\
\(\Re,\Im\) & Real- og imaginærdel. & -- & \S\ref{sec:bode}, \S\ref{sec:pilead}\\
\(\angle(\cdot)\) & Fase/vinkel af kompleks størrelse. & grader & \S\ref{sec:bode}, \S\ref{sec:pilead}\\
\(\operatorname{atan2}\) & Kvadrantkorrekt vinkelfunktion. & Funktion & \S\ref{sec:bode}, \S\ref{sec:pilead}, \S\ref{sec:python}\\
\(\arctan\), \(\sin\), \(\sin^{-1}\), \(\sqrt{\cdot}\), \(\log_{10}\), \(e^{(\cdot)}\) & Matematiske standardfunktioner brugt i fase-, lead-, Bode- og tidsresponsformler. & Funktioner & \S\ref{sec:time}, \S\ref{sec:bode}, \S\ref{sec:pilead}, \S\ref{sec:limits}\\
\(\lim\), \(\max\), \(\min\) & Grænseværdi og maksimum/minimum; bruges til DC-gain, final value og actuatorgrænser. & Operatorer & \S\ref{sec:model}, \S\ref{sec:feedback}, \S\ref{sec:limits}\\
\(\pi\), \(\infty\), \(\%\), \(^\circ\) & Matematisk konstant, uendelighed, procent og grader. & Konstanter/enheder & \S\ref{sec:start}, \S\ref{sec:time}, \S\ref{sec:bode}, \S\ref{sec:pilead}\\
\(<,>,\leq,\geq\) & Ulighedstegn til stabilitetsintervaller, parameterkrav og actuatorlimits. & Relationer & \S\ref{sec:start}, \S\ref{sec:model}, \S\ref{sec:limits}, \S\ref{sec:python}\\
\(\sum F\), \(m\), \(\ddot x\) & Newtons 2. lov i mekanisk modellering: kraftsum, masse og acceleration. & N, kg, m/s\(^2\) & \S\ref{sec:model}\\
\end{longtable}
\normalsize
